% =====================================================================
%  Lecture 14 -- 5G Networks: Introduction                (English edition)
% =====================================================================
\documentclass[11pt,a4paper]{article}
\providecommand\ACNroot{../..}
\input{\ACNroot/templates/lecture-en}

\ACNset{number}{14}
\ACNset{title}{5G Networks: Introduction}
\ACNset{shorttitle}{5G introduction}
\ACNset{titleother}{مقدمه‌ای بر شبکه‌های نسل پنجم}
\ACNset{part}{Part IV -- 5G system and procedures (Lectures 14, 19--22)}
\ACNset{time}{6--8 hours (3 h reading, 3--5 h exercises)}
\ACNset{level}{Intermediate}
\ACNset{prereq}{Lectures 5--7 and 9; OFDM basics; Shannon capacity}
\ACNset{gen}{5}

\begin{document}
\ACNcover

% ---------------------------------------------------------------------
\section{Learning objectives}
After studying this lecture you should be able to:
\begin{objectives}
  \item \textbf{State} the IMT-2020 usage scenarios and key minimum requirements, and \textbf{relate}
        each to a design choice of 5G.
  \item \textbf{Explain} NR's scalable numerology and \textbf{compute} slot and symbol durations for
        each subcarrier spacing.
  \item \textbf{Compute} the approximate peak data rate of an NR carrier with the 3GPP formula.
  \item \textbf{Compare} non-standalone (EN-DC) and standalone deployments.
  \item \textbf{Estimate} latency, connection density and area capacity for given deployments, and
        \textbf{judge} whether a requirement is met.
  \item \textbf{Outline} the 5G core and the topics of the following lectures.
\end{objectives}

% ---------------------------------------------------------------------
\section{Prerequisite knowledge}
You need the LTE architecture and radio protocols (Lectures~5--7), the EPC (Lecture~9), the structure of
an OFDM resource grid (subcarriers, symbols, cyclic prefix, resource blocks) and Shannon's capacity
formula $C = B\log_2(1+\SNR)$ \cite{shannon1948mathematical}.

\paragraph{Key terms} The vocabulary of this lecture with its Persian equivalents:
\begin{keyterms}
IMT-2020 & \fa{IMT-2020} & ITU-R's name for the 5G family of radio interfaces and its requirements. \\
eMBB / URLLC / mMTC & \fa{سه سناریوی استفاده} & Enhanced broadband, ultra-reliable low latency, massive machine type. \\
NR & \fa{رادیوی جدید} & 3GPP's 5G radio interface (Release~15 onward). \\
Numerology $\mu$ & \fa{عددشناسی} & Subcarrier spacing $15\cdot2^\mu$\,kHz and the matching slot length. \\
FR1 / FR2 & \fa{بازه‌های بسامدی ۱ و ۲} & Sub-7\,GHz and millimetre-wave frequency ranges. \\
Bandwidth part & \fa{بخش پهنای باند} & Subset of a carrier configured for a UE. \\
EN-DC & \fa{اتصال دوگانهٔ E-UTRA--NR} & LTE master and NR secondary node with the EPC (NSA). \\
SA / NSA & \fa{مستقل / غیرمستقل} & NR with the 5G core / NR anchored on LTE. \\
5GC & \fa{هستهٔ نسل پنجم} & Service-based core network (AMF, SMF, UPF, \ldots). \\
Network slice & \fa{برش شبکه} & Logical end-to-end network for a service type. \\
\end{keyterms}

% ---------------------------------------------------------------------
\section{Introduction and motivation}
LTE was designed for one dominant use: mobile broadband to smartphones. By the mid-2010s operators and
industries expected three different kinds of demand: far higher capacity for video and fixed wireless
access, wireless control of machines with reliability and latency comparable to wired fieldbuses, and
connectivity for very large numbers of cheap, battery-powered sensors. ITU-R captured these as the
three IMT-2020 usage scenarios (\cref{fig:usage}) \cite{itu2015m2083} and set minimum technical
requirements for any candidate radio interface \cite{itu2017m2410}. 3GPP's answer is the 5G system: the
NR radio interface \cite{ts38300,dahlman2020nr} and the service-based 5G core \cite{ts23501}.

This lecture opens Part~IV. It surveys the requirements, the main ideas of NR that address them, and the
migration from LTE; the core network, procedures, slicing and QoS follow in Lectures~18--22, while
software-defined networking and virtualisation, the technologies on which the 5G core is built, are the
topic of Part~III.

\begin{acnfigure}
  \acnfig{lecture-14/usage-scenarios}
  \caption{The three IMT-2020 usage scenarios with example applications \cite{itu2015m2083}.}
  \label{fig:usage}
\end{acnfigure}

% ---------------------------------------------------------------------
\section{Fundamental concepts}

\subsection{Requirements}
\begin{table}[htbp]
\centering\small
\caption{Selected IMT-2020 minimum technical requirements (ITU-R M.2410) \cite{itu2017m2410}.}\label{tab:req}
\rowcolors{2}{acnmist}{white}
\begin{tabularx}{\linewidth}{@{}l l Y@{}}
\toprule
\rowcolor{white}Requirement & Value & Scenario / remark \\
\midrule
Peak data rate             & DL 20\,Gbit/s, UL 10\,Gbit/s  & eMBB \\
Peak spectral efficiency   & DL 30, UL 15 bit/s/Hz         & eMBB \\
User-experienced data rate & DL 100\,Mbit/s, UL 50\,Mbit/s & eMBB, dense urban (5th percentile) \\
Area traffic capacity      & 10\,Mbit/s/m$^2$              & eMBB, indoor hotspot, DL \\
User-plane latency         & 4\,ms (eMBB), 1\,ms (URLLC)   & One way, unloaded \\
Control-plane latency      & 20\,ms                        & Idle to active \\
Connection density         & $10^6$ devices/km$^2$         & mMTC \\
Reliability                & $1-10^{-5}$ for 32 bytes in 1\,ms & URLLC \\
Mobility                   & up to 500\,km/h               & High-speed trains \\
Bandwidth                  & at least 100\,MHz; up to 1\,GHz above 6\,GHz & -- \\
\bottomrule
\end{tabularx}
\end{table}

No single deployment has to meet all values at once: each requirement is evaluated in its own test
environment. \Cref{fig:capabilities} shows the improvement over IMT-Advanced (4G) that ITU-R envisaged for
each key capability.

\begin{acnfigure}
  \acnfig{lecture-14/capabilities}
  \caption{Improvement factors of IMT-2020 key capabilities over IMT-Advanced, as envisaged in ITU-R M.2083
           \cite{itu2015m2083}. Latency is shown as its inverse.}
  \label{fig:capabilities}
\end{acnfigure}

\subsection{Design principles of NR}
\begin{itemize}
  \item \textbf{Wide spectrum range}: FR1 (410\,MHz--7.125\,GHz) and FR2 (24.25--52.6\,GHz, extended in later
        releases), with carriers up to 100\,MHz in FR1 and 400\,MHz in FR2 \cite{ts38300}.
  \item \textbf{Scalable numerology}: subcarrier spacing $15\cdot2^\mu$\,kHz (\cref{sec:num}).
  \item \textbf{Ultra-lean design}: no always-on cell-specific reference signals; signals are sent only when
        needed, which saves energy and reduces interference.
  \item \textbf{Beam-centric operation}: synchronisation signal blocks are swept in beams; beam management is
        built in, which is essential at millimetre waves.
  \item \textbf{Flexible scheduling}: slots and mini-slots (2, 4 or 7 symbols), dynamic TDD, bandwidth parts.
  \item \textbf{New channel coding}: LDPC for data, polar codes for control \cite{ts38211}.
  \item \textbf{Forward compatibility}: reserved resources so that future features do not break legacy UEs.
\end{itemize}

\subsection{Deployment options}
3GPP defined several ways to combine LTE, NR, the EPC and the 5GC. The two that matter in practice
(\cref{fig:options}):
\begin{itemize}
  \item \textbf{Non-standalone (NSA), option 3 family (EN-DC)}: the UE is connected to an LTE eNB (master node,
        control plane to the MME) and an NR en-gNB (secondary node, adding capacity) \cite{ts37340}. It reuses the
        EPC and allowed the first 5G launches in 2019 without a new core, but provides none of the new core
        features (slicing, new QoS, service-based architecture).
  \item \textbf{Standalone (SA), option 2}: the gNB connects to the 5GC over N2 (control) and N3 (user plane).
        All 5G features, including low latency and slicing, require SA.
\end{itemize}

\begin{acnfigure}
  \resizebox{\linewidth}{!}{\acnfig{lecture-14/deployment-options}}
  \caption{Non-standalone EN-DC (option 3x) with the EPC, and standalone NR (option 2) with the 5GC.}
  \label{fig:options}
\end{acnfigure}

% ---------------------------------------------------------------------
\section{Detailed technical explanation}

\subsection{Frame structure}
NR keeps LTE's 10\,ms frame and 1\,ms subframe. A slot has 14 OFDM symbols (normal cyclic prefix); with
numerology $\mu$ a subframe contains $2^\mu$ slots (\cref{fig:numerology}). A resource block is 12 subcarriers
in frequency; a carrier has at most 275 resource blocks. Because 1\,ms is a fixed grid for all numerologies,
symbols of different numerologies align at subframe boundaries, which allows mixing them on one carrier
\cite{ts38211}.

\begin{acnfigure}
  \acnfig{lecture-14/numerology}
  \caption{Slots per 1\,ms subframe for numerologies $\mu = 0$--3.}
  \label{fig:numerology}
\end{acnfigure}

\subsection{Initial access and beams}
A gNB transmits synchronisation signal blocks (SSBs: PSS, SSS and PBCH) in a burst, each in a different beam.
The UE detects the best SSB, reads the system information and performs random access on the RACH occasion
associated with that SSB; the gNB thereby learns the best beam. During a connection, beams are refined with
CSI-RS measurements and reports. In FR1 a few wide beams suffice; in FR2 up to 64 narrow beams compensate the
higher path loss with array gain.

\subsection{From EPC to 5GC}
The 5G core separates access and mobility management (AMF) from session management (SMF) and the user plane
(UPF), stores subscription data in the UDM/UDR, authenticates through the AUSF and controls policy in the PCF.
Network functions expose services over HTTP/2 and discover each other through the NRF \cite{ts23501}. The QoS
model uses QoS flows within PDU sessions instead of EPS bearers, and network slices allow several logical
networks on one infrastructure. These topics are developed in Lectures~18--22.

% ---------------------------------------------------------------------
\section{Mathematical models}

\subsection{Numerology}\label{sec:num}
\paragraph{Definitions} For $\mu = 0,1,2,3,4$:
\begin{equation}\label{eq:numerology}
  \Delta f = 15\cdot2^{\mu}\ \text{kHz}, \qquad
  T_{\mathrm{slot}} = \frac{1\ \text{ms}}{2^{\mu}}, \qquad
  \bar T_{\mathrm{sym}} = \frac{T_{\mathrm{slot}}}{14},
\end{equation}
where $\bar T_{\mathrm{sym}}$ is the mean symbol duration including the cyclic prefix. The normal cyclic prefix is
about $4.7/2^\mu$\,\textmu s.
\paragraph{Interpretation} Wider spacing shortens symbols and slots (lower latency) and tolerates more phase
noise and Doppler (needed in FR2), but its shorter cyclic prefix tolerates less delay spread, which suits small
cells. FR1 macro cells typically use 15 or 30\,kHz; FR2 uses 120\,kHz.

\subsection{Peak data rate}\label{sec:peak}
\paragraph{Formula} 3GPP approximates the maximum data rate of a carrier as \cite{ts38306}
\begin{equation}\label{eq:peak}
  R \approx 10^{-6}\, v\, Q_m\, f\, R_{\max}\,
          \frac{N_{\mathrm{PRB}}\cdot 12}{T_s^{\mu}}\,\bigl(1-\mathrm{OH}\bigr)\ \text{Mbit/s},
  \qquad T_s^{\mu} = \frac{10^{-3}}{14\cdot2^{\mu}}\ \text{s},
\end{equation}
with $v$ MIMO layers, $Q_m$ bits per modulation symbol, scaling factor $f\le1$, maximum code rate
$R_{\max} = 948/1024$, $N_{\mathrm{PRB}}$ resource blocks and overhead OH (0.14 in FR1 DL, 0.18 in FR2 DL, 0.08 in FR1
UL, 0.10 in FR2 UL). For carrier aggregation the rates of the carriers are added.
\paragraph{Example} FR1, 100\,MHz at 30\,kHz ($N_{\mathrm{PRB}} = 273$), 4 layers, 256-QAM ($Q_m = 8$), $f = 1$:
$R \approx 2.34$\,Gbit/s (\cref{ex:peak}).

\subsection{User-plane latency}\label{sec:lat}
\paragraph{Model} A simple estimate of one-way radio latency is
\begin{equation}\label{eq:latency}
  T_{\mathrm{UP}} \approx T_{\mathrm{align}} + T_{\mathrm{tx}} + T_{\mathrm{proc}}
  \approx \tfrac12 T_{\mathrm{slot}} + T_{\mathrm{slot}} + 2\,T_{\mathrm{slot}},
\end{equation}
counting on average half a slot to wait for the next transmission opportunity, one slot of transmission and
about two slots of processing at transmitter and receiver.
\paragraph{Example} $\mu = 0$: 3.5\,ms; $\mu = 1$: 1.75\,ms; $\mu = 3$: 0.44\,ms. Mini-slots and pre-configured
(grant-free) uplink shorten $T_{\mathrm{align}}$ and $T_{\mathrm{tx}}$ further.
\paragraph{Limitations} The model ignores scheduling requests, HARQ retransmissions, queueing and the transport
network and core; it shows why the URLLC target of 1\,ms requires larger numerologies or mini-slots.

\subsection{Connection density and area capacity}\label{sec:density}
\paragraph{Models} With device density $\rho$ (devices/km$^2$) and cell radius $r$, a cell serves
$N = \rho\pi r^2$ devices; if each accesses the network once per period $T$, the access rate is $N/T$. For area
capacity, a cell with bandwidth $B$ and mean spectral efficiency $\bar\eta$ covering area $A_{\mathrm{cell}}$ offers
\begin{equation}\label{eq:area}
  C_{\mathrm{area}} = \frac{B\,\bar\eta}{A_{\mathrm{cell}}}.
\end{equation}
\paragraph{Examples} $\rho = 10^6$/km$^2$, $r = 500$\,m: $N = 7.9\times10^5$ devices; one report every 2\,h gives 109
accesses/s per cell (\cref{ex:density}). For $C_{\mathrm{area}} = 10$\,Mbit/s/m$^2$ with $B = 400$\,MHz and $\bar\eta =
5$\,bit/s/Hz, each cell may cover only $A_{\mathrm{cell}} = 200$\,m$^2$: area capacity is achieved by densification,
not by spectral efficiency alone.

% ---------------------------------------------------------------------
\section{Network architecture and diagrams}
\Cref{fig:options} summarises the two main architectures. In NSA the EPC remains the core, and the NR node is an
extension of the LTE radio network (it has no control-plane connection to the core). In SA, the gNB (split into
central and distributed units, Lecture~23) connects to the 5GC; LTE and NR then interwork through the core
(N26 between AMF and MME) for mobility and voice fallback. Most operators started with NSA and introduced SA
later, often first for specific segments such as fixed wireless access or enterprise slices.

% ---------------------------------------------------------------------
\section{Step-by-step procedures}

\begin{procedurebox}{Adding an NR secondary node in EN-DC}\label{proc:endc}
\begin{steps}
  \item The UE is connected to the LTE eNB (master node) and reports NR measurements configured by it.
  \item The eNB decides to add an NR leg and sends SgNB Addition Request over X2 to the en-gNB.
  \item The en-gNB allocates resources and answers with the NR radio configuration; the eNB forwards it in
        RRC Connection Reconfiguration.
  \item The UE performs random access on the NR cell and confirms; the eNB informs the MME (E-RAB modification)
        so that the S-GW can send the user plane directly to the en-gNB (option 3x).
  \item Data now flow over NR (and LTE, if the bearer is split); mobility and control remain with LTE.
\end{steps}
\end{procedurebox}

\begin{procedurebox}{Initial access in standalone NR}\label{proc:access}
\begin{steps}
  \item Search the SSB raster, detect PSS/SSS, determine the cell ID and the best SSB beam.
  \item Read the MIB on PBCH and SIB1 (scheduled by it) to obtain RACH configuration and cell parameters.
  \item Transmit a preamble in the RACH occasion of the best SSB; receive the random-access response (Lecture~8).
  \item Complete RRC set-up and send the Registration Request to the AMF (Lecture~20).
\end{steps}
\end{procedurebox}

% ---------------------------------------------------------------------
\section{Worked examples}

\begin{examplebox}{Peak rate of an FR1 carrier}\label{ex:peak}
\textbf{Problem.} Evaluate \cref{eq:peak} for 100\,MHz at 30\,kHz, 4 layers, 256-QAM, $f = 1$, DL.

\textbf{Solution.} $T_s^{1} = 10^{-3}/28 = 35.7$\,\textmu s. $N_{\mathrm{PRB}}\cdot12/T_s = 273\times12/35.7\times10^{-6} = 9.17\times10^{7}$
REs/s. $R = 4\times8\times0.926\times9.17\times10^{7}\times0.86 = 2.34\times10^{9}$\,bit/s. To reach 20\,Gbit/s, NR needs
FR2 bandwidths, more layers and carrier aggregation.
\end{examplebox}

\begin{examplebox}{Latency and numerology}\label{ex:lat}
\textbf{Problem.} Which numerologies satisfy the 1\,ms URLLC target with the model of \cref{eq:latency}?

\textbf{Solution.} $T_{\mathrm{UP}} = 3.5\,T_{\mathrm{slot}} \le 1$\,ms requires $T_{\mathrm{slot}} \le 0.286$\,ms, i.e.\
$\mu \ge 2$ (0.875\,ms at 60\,kHz). With 30\,kHz, mini-slots of 2 symbols reduce the transmission and alignment terms and can meet
the target as well.
\end{examplebox}

\begin{examplebox}{mMTC access load}\label{ex:density}
\textbf{Problem.} A cell of radius 500\,m in an area with $10^6$ devices/km$^2$; each device sends one report every two
hours. Compute the access rate and compare with LTE's typical RACH capacity of the order of a few hundred
attempts per second per cell.

\textbf{Solution.} $N = 10^6\times\pi\times0.25 = 7.85\times10^5$; rate $= 7.85\times10^5/7200 = 109$\,s$^{-1}$ on average --
manageable, but synchronised reporting (all devices at the full hour) would be catastrophic (Lecture~8). mMTC
therefore relies on randomised schedules, early data transmission and NB-IoT/LTE-M features, not only on raw
capacity.
\end{examplebox}

% ---------------------------------------------------------------------
\section{Technology comparison}
\begin{table}[htbp]
\centering\small
\caption{LTE and NR compared.}\label{tab:cmp}
\rowcolors{2}{acnmist}{white}
\begin{tabularx}{\linewidth}{@{}l Y Y@{}}
\toprule
\rowcolor{white} & LTE (4G) & NR (5G) \\
\midrule
Spectrum          & Below 6\,GHz             & FR1 and FR2 (mmWave) \\
Max carrier bandwidth & 20\,MHz (aggregation up to 32 carriers) & 100\,MHz (FR1), 400\,MHz (FR2) \\
Subcarrier spacing & 15\,kHz (fixed)          & $15\cdot2^\mu$\,kHz \\
Scheduling unit   & 1\,ms subframe           & Slot ($1/2^\mu$\,ms), mini-slot \\
Reference signals & Always-on CRS            & On demand (ultra-lean) \\
Channel coding    & Turbo (data), convolutional (control) & LDPC (data), polar (control) \\
Beamforming       & Added in later releases  & Native, including initial access \\
Core              & EPC                      & EPC (NSA) or 5GC (SA) \\
QoS unit          & EPS bearer               & QoS flow \\
\bottomrule
\end{tabularx}
\end{table}

\paragraph{Evolution, not replacement} LTE continues to evolve (LTE-M, NB-IoT, and dynamic spectrum sharing with
NR on the same carrier), and NR reuses many LTE concepts: OFDM, HARQ, the layer-2 stack, and the
attach/registration and bearer/flow logic of the core.

% ---------------------------------------------------------------------
\section{Key takeaways}
\begin{takeawaybox}
\begin{itemize}
  \item IMT-2020 defines three scenarios (eMBB, URLLC, mMTC) and minimum requirements such as 20\,Gbit/s peak,
        1\,ms URLLC latency and $10^6$ devices/km$^2$.
  \item NR: wide spectrum (FR1, FR2), numerology $15\cdot2^\mu$\,kHz, slots of $1/2^\mu$\,ms, lean and beam-based.
  \item Peak rate $\propto v\,Q_m\,R_{\max}\,N_{\mathrm{PRB}}\cdot12/T_s\,(1-\mathrm{OH})$; 100\,MHz FR1 gives about 2.3\,Gbit/s with 4 layers.
  \item NSA (EN-DC) reuses the EPC; SA with the 5GC unlocks slicing, new QoS and low latency.
  \item Area capacity comes from densification; latency from numerology and mini-slots; density from access design.
\end{itemize}
\end{takeawaybox}

% ---------------------------------------------------------------------
\section{Exercises}

\begin{exercisebox}{1}{Which scenario?}\label{exr:scenario}
Assign each service to eMBB, URLLC or mMTC and name its most demanding requirement: (a) 8K video streaming;
(b) motion control in a factory; (c) water meters; (d) remote surgery; (e) fixed wireless access to homes.
\end{exercisebox}

\begin{exercisebox}{1}{Why NSA first?}\label{exr:nsa}
Give two reasons why most operators launched 5G as NSA, and one capability they could not offer until SA.
\end{exercisebox}

\begin{exercisebox}{2}{FR2 peak rate}\label{exr:fr2}
Compute the peak DL rate of a 400\,MHz FR2 carrier at 120\,kHz ($N_{\mathrm{PRB}} = 264$) with 2 layers and 64-QAM.
How many such carriers would be needed for 20\,Gbit/s?
\end{exercisebox}

\begin{exercisebox}{2}{Cyclic prefix and cell size}\label{exr:cp}
The normal cyclic prefix is about $4.7/2^\mu$\,\textmu s. Compute it for $\mu = 0$--3 and the corresponding path-length
difference it can absorb. Which numerologies suit a rural macro cell with a delay spread of 2\,\textmu s?
\end{exercisebox}

\begin{exercisebox}{3}{Peak spectral efficiency}\label{exr:se}
Using \cref{eq:peak}, compute the peak spectral efficiency of a 100\,MHz FR1 carrier at 30\,kHz with 8 DL layers and
256-QAM. Does it meet the IMT-2020 requirement of 30\,bit/s/Hz? Which assumptions make this a theoretical
value?
\end{exercisebox}

\begin{exercisebox}{3}{Densification}\label{exr:dense}
An indoor hotspot must provide 10\,Mbit/s/m$^2$. Compare the cell area needed with (a) 400\,MHz at 5\,bit/s/Hz and
(b) 100\,MHz at 7\,bit/s/Hz. Discuss the implications for backhaul and for the core network.
\end{exercisebox}

\subsection*{Solutions}
\addcontentsline{toc}{subsection}{Solutions}

\begin{solutionbox}{\ref{exr:scenario}}
(a) eMBB -- user data rate; (b) URLLC -- latency and reliability; (c) mMTC -- connection density and battery life;
(d) URLLC (with eMBB video) -- reliability and latency; (e) eMBB -- capacity per area and per user.
\end{solutionbox}

\begin{solutionbox}{\ref{exr:nsa}}
NSA reused the existing EPC (no new core, faster time to market) and LTE coverage for control and mobility, while
NR added capacity; early devices and spectrum were also tuned for it. Features needing the 5GC -- network slicing,
QoS flows, service-based exposure, low-latency edge user planes, and NR-only coverage -- required SA.
\end{solutionbox}

\begin{solutionbox}{\ref{exr:fr2}}
$T_s^{3} = 10^{-3}/112 = 8.93$\,\textmu s; $R = 2\times6\times0.926\times264\times12/8.93\times10^{-6}\times0.82 = 3.23$\,Gbit/s. For
20\,Gbit/s about $20/3.23 = 6.2$, i.e.\ 7 carriers -- or more layers and higher-order modulation.
\end{solutionbox}

\begin{solutionbox}{\ref{exr:cp}}
$\mu = 0$: 4.69\,\textmu s (1406\,m); $\mu = 1$: 2.34\,\textmu s (703\,m); $\mu = 2$: 1.17\,\textmu s (352\,m); $\mu = 3$: 0.59\,\textmu s
(176\,m), using $d = c\,T_{\mathrm{CP}}$. With 2\,\textmu s delay spread, $\mu = 0$ (and marginally $\mu = 1$) are suitable;
60\,kHz needs the extended CP, and 120\,kHz is for small FR2 cells.
\end{solutionbox}

\begin{solutionbox}{\ref{exr:se}}
With 8 layers the rate doubles to $4.67$\,Gbit/s, i.e.\ $46.7$\,bit/s/Hz over 100\,MHz -- above 30. The value assumes
8 spatially separable layers, 256-QAM at the highest code rate (very high SINR), the fixed overhead factor and full
allocation to one UE; real systems reach a fraction of it.
\end{solutionbox}

\begin{solutionbox}{\ref{exr:dense}}
(a) $400\times10^6\times5 = 2$\,Gbit/s per cell, so $A = 2\times10^9/10^7 = 200$\,m$^2$; (b) $0.7$\,Gbit/s, so 70\,m$^2$. Either way
cells are rooms, not streets. Each cell needs multi-gigabit backhaul (fibre), and the core must handle traffic of
$10$\,Tbit/s per km$^2$ of such hotspots locally -- an argument for edge user planes (Lecture~23).
\end{solutionbox}

% ---------------------------------------------------------------------
\section{Further reading}
ITU-R M.2083 \cite{itu2015m2083} sets out the vision and M.2410 \cite{itu2017m2410} the requirements; 3GPP's own scenarios
are in TR~38.913 \cite{tr38913} and service requirements in TS~22.261 \cite{ts22261}. NR is described in TS~38.300
\cite{ts38300} and TS~38.211 \cite{ts38211}, EN-DC in TS~37.340 \cite{ts37340} and the 5G system in TS~23.501
\cite{ts23501}. Dahlman et al.\ \cite{dahlman2020nr} is the standard textbook on NR.

\printbibliography[heading=subbibliography,title={\ACNstrRefs}]
\end{document}
